\section{Evaluation}
\label{sec:eval}

We organize the evaluation around seven research questions, preceded by a validation step that
establishes the simulator and the proof-of-concept are sound. RQ1--RQ5 exercise the model; RQ6
cross-validates it against a discrete-event testbed with real signatures; RQ7 applies it to the
documented parameters of deployed systems. All numbers are produced by the artifact
(\S\ref{sec:impl}); figure and table references are to artifact output.

\paragraph{Validation and proof-of-concept (RQ0).}
The Monte-Carlo BFT harness reproduces the closed-form $\psucc$ of Eq.~\eqref{eq:psucc} across all
tested $(N,\varphi)$ with maximum absolute deviation $\ArtifactMaxValErr$ (tolerance $0.02$),
validating both the harness and the closed form. A sub-floor purchase ($k<F{+}1$) yields a conflicting
certificate in \emph{zero} of every trial, confirming Lemma~\ref{lem:intersection} empirically. The
end-to-end PoC, on a committee of $N=\ArtifactPoCN$ ($F=10$, floor $k=\ArtifactPoCk$), drives both
branches to the quorum $Q=21$ under an engineered split, emits $\ArtifactPoCEvidence$ valid
equivocation proofs, and the \textsc{PayToEquivocate} escrow pays exactly those $\ArtifactPoCBribes$
and rejects a forged (non-conflicting) proof. The attack is therefore realizable end-to-end in the
local model, not merely asserted.

\paragraph{The floor is necessary, not sufficient.}
Figure~\ref{fig:psucc} plots $\psucc$ against the view-split $\varphi$ for $k\in\{F{+}1,F{+}2,F{+}3\}$,
with the Monte-Carlo overlay matching the closed form. At the floor $k=F{+}1$, success requires an
exact even split, so $\psucc$ peaks at only $\ArtifactPsuccFloor$ at $\varphi=\tfrac12$ and falls off
sharply for asymmetric splits; each extra purchased equivocator widens the feasible band
(Eq.~\eqref{eq:band}) and raises $\psucc$. This quantifies the central caveat: buying the $F{+}1$
floor is necessary but, without engineering a balanced honest split, far from sufficient.

\begin{figure}[t]
\centering
\includegraphics[width=\linewidth]{fig_psucc.pdf}
\caption{$\psucc$ vs.\ view-split $\varphi$ for the floor and two larger purchases ($N=31$). Markers:
Monte-Carlo harness; lines: closed form (Eq.~\eqref{eq:psucc}).}
\label{fig:psucc}
\end{figure}

\paragraph{RQ1: When does the settlement window open?}
The window is nearly a step function---there is almost no middle ground between closed and open. The
window-open probability $\pwin$ crosses $0.5$ within $1\%$ of $\rho=\Tacc/\Tsettle=1$ and saturates
toward $1$ for $\rho\gg1$ (Figure~\ref{fig:pwin}), with higher settle-time variance softening the
transition, exactly as the model predicts. The finality-gated baseline (Proposition~\ref{prop:immunity})
is identically zero for all $\rho$---the flat line at the bottom---so closing the window is not a matter
of degree but of construction.

\begin{figure}[t]
\centering
\includegraphics[width=\linewidth]{fig_pwin.pdf}
\caption{RQ1. Window-open probability $\pwin$ vs.\ latency ratio $\rho$, for two settle-time
coefficients of variation, plus the finality-gated immunity baseline ($\pwin\equiv0$).}
\label{fig:pwin}
\end{figure}

\paragraph{RQ2: How does the bribe cost scale?}
The bribe is strikingly small relative to the committee: at $N=31,\,q=0.5,\,R=0.5s$ the minimum total
bribe is just $8.25s$---the adversary pays $F{+}1{=}11$ validators a \emph{fraction} $q$ of their
stake-at-risk, not the whole committee's stake. Figure~\ref{fig:cost} plots $B=(F{+}1)q(s{+}R)$ against
committee size $N$ for several $q$: cost grows only \emph{linearly} in $N$ (through $F{+}1$) and
linearly in $q$. The policy reading is that raising $N$ increases the (G2) cost only linearly,
whereas closing (G1) removes the threat outright---consistent with Proposition~\ref{prop:immunity} and
quantified in RQ5.

\begin{figure}[t]
\centering
\includegraphics[width=\linewidth]{fig4_cost.pdf}
\caption{RQ2. Minimum total bribe $B$ vs.\ committee size $N$ for $q\in\{0.25,0.5,0.75,1.0\}$,
$R=0.5s$.}
\label{fig:cost}
\end{figure}

\paragraph{RQ3: When is the attack profitable?}
Profit lives in a single quadrant: both gates must hold at once. Figure~\ref{fig:phase} shows
$\mathbb{E}[\Pi]$ over the $(\rho,\Eval)$ plane is positive only in the upper-right region where
$\rho>1$ \emph{and} $\Eval>\Eval^{*}$; the finality-gated baseline collapses the entire plane to
non-positive profit. Table~\ref{tab:breakeven} reports the break-even value
$\Eval^{*}=(F{+}1)q(s{+}R)+\Cop$ across $N$ and $q$. Figure~\ref{fig:profitqV}
slices the profitable region over $(q,\Eval)$ at a window-open $\rho$, showing the linear shift of the
break-even frontier with $q$. Sensitivity to the adversary's split-control reliability is direct: the
break-even value scales as $1/(\psucc\pwin)$, so a weak adversary ($\psucc=0.2$) needs roughly
$4{\times}$ the value of a strong one ($\psucc=0.85$) to clear (G2). We do not treat $\psucc=0.85$ as a
default: it is the view-control regime; an uncontrolled split caps $\psucc$ at $\ArtifactPsuccFloor$ at
the $F{+}1$ floor, and reaching $0.85$ by purchase alone costs $k\approx F{+}6$ at a $1.45\times$ bribe
(\S\ref{sec:floor}). The phase diagram (Fig.~\ref{fig:phase}) is thus the strong-adversary upper edge of
a band whose lower edge is the floor.

\begin{table}[t]
\centering
\caption{RQ3. Break-even extractable value $\Eval^{*}=(F{+}1)q(s{+}R)+\Cop$ (units of $s$; $R=0.5s$,
$\Cop=1$).}
\label{tab:breakeven}
\small
% AUTO-GENERATED by artifact/make_tables.py -- do not edit by hand.
\begin{tabular}{@{}lcccc@{}}
\toprule
$N$ ($k{=}F{+}1$) & $q{=}0.25$ & $q{=}0.5$ & $q{=}0.75$ & $q{=}1.0$ \\
\midrule
$4$ ($2$) & 1.75 & 2.50 & 3.25 & 4.00 \\
$7$ ($3$) & 2.12 & 3.25 & 4.38 & 5.50 \\
$16$ ($6$) & 3.25 & 5.50 & 7.75 & 10.00 \\
$31$ ($11$) & 5.12 & 9.25 & 13.38 & 17.50 \\
$64$ ($22$) & 9.25 & 17.50 & 25.75 & 34.00 \\
$100$ ($34$) & 13.75 & 26.50 & 39.25 & 52.00 \\
\bottomrule
\end{tabular}

\end{table}

\begin{figure}[t]
\centering
\includegraphics[width=\linewidth]{fig3_phase.pdf}
\caption{RQ3. Profit phase diagram $\mathbb{E}[\Pi]$ over $(\rho,\Eval)$ ($N=31$, $q=0.5$,
$\psucc=0.85$; uncontrolled-split floor $\ArtifactPsuccFloor$). The $\mathbb{E}[\Pi]=0$ contour bounds
the profitable region (both gates hold).}
\label{fig:phase}
\end{figure}

\begin{figure}[t]
\centering
\includegraphics[width=\linewidth]{fig_profit_qV.pdf}
\caption{RQ3. Profit over $(q,\Eval)$ at a window-open latency ratio.}
\label{fig:profitqV}
\end{figure}

\paragraph{RQ4: How does reconfiguration change the race?}
Figure~\ref{fig:reconf-profit} (\S\ref{sec:reconf}) reports $\mathbb{E}[\Pi(a)]$ across the epoch for
the six handoff designs. The hard-cutoff design is profitable mid-epoch but collapses to
$-B-\Cop$ as $a\to\tau$ (residual time $r\to0$): last-block attacks are weak. The unsafe-carryover and
reconfiguration-delay designs instead \emph{rise} toward the boundary---old-epoch artifacts survive,
$\Hreconf(a)$ widens the deadline, and $q(a)$ falls (cheaper bribes)---confirming
Proposition~\ref{prop:reconf-residual}. The accountability-gated carryover stays flat and negative
($\pwin=0$, Proposition~\ref{prop:reconf-safe}), and rolling reconfiguration shows a smooth optimum
with no last-block cliff (Counterexample~\ref{cex:rolling}).

\paragraph{RQ5: Which defenses close the attack?}
Only the two levers that touch the window flip the sign of profit; the two that only raise the cost
leave the attack firmly profitable (Table~\ref{tab:rq5defenses}). From a profitable baseline
($\mathbb{E}[\Pi]=+38.9$), finality-gated receipts drive profit to $-9.2$ (the immunity baseline,
$\pwin\!\to\!0$), and faster evidence---which lowers $\Tacc$, hence both $\pwin$ and (via larger $q$)
the bribe---drives it to $-0.5$. Purely economic levers do not: persistent slashability (raising $q$,
hence $B$) leaves profit at $+32.3$, and a larger committee (raising $B$ linearly) at $+21.7$. Both
the table and Figure~\ref{fig:defenses} are the quantitative form of the master-variable claim:
closing the window dominates raising the cost.

\begin{table}[t]
\centering
\caption{RQ5. Each defensive lever from a profitable baseline ($N=31$, $V=60s$, $\psucc=0.85$): the
gate it attacks (G1 $=$ window $\Tsettle<\Tacc$; G2 $=$ cost $V>B$), the resulting window probability
$\pwin$, bribe $B$ (units of $s$), expected profit $\mathbb{E}[\Pi]$, and whether it makes the attack
unprofitable.}
\label{tab:rq5defenses}
\small
\setlength{\tabcolsep}{4.5pt}
% AUTO-GENERATED by artifact/make_tables.py -- do not edit by hand.
\begin{tabular}{@{}l c r r r c@{}}
\toprule
Defensive lever & Gate & $\pwin$ & $B$ & $\mathbb{E}[\Pi]$ & Closes? \\
\midrule
None (baseline) & -- & 0.95 & 8.25 & +39.3 & \nomark \\
Finality-gating & G1 & 0.00 & 8.25 & -9.2 & \yes \\
Persistent slash. & G2 & 0.95 & 14.85 & +32.7 & \nomark \\
Fast evidence & G1+G2 & 0.30 & 14.85 & -0.4 & \yes \\
Larger committee & G2 & 0.95 & 25.50 & +22.0 & \nomark \\
\bottomrule
\end{tabular}

\end{table}

\begin{figure}[t]
\centering
\includegraphics[width=\linewidth]{fig_defenses.pdf}
\caption{RQ5. Effect of each defensive lever on $\mathbb{E}[\Pi]$ ($N=31$, $V=60s$).}
\label{fig:defenses}
\end{figure}

\paragraph{RQ6: Does a running protocol reproduce the model?}
The analyses above are validated against an \emph{independent mechanism}. The discrete-event testbed
(\S\ref{sec:impl}) runs an actual BFT commit round over a modeled network with real Ed25519
signatures and derives $\Tacc$ and $\Tsettle$ from event timing rather than sampling them from the
model. Sweeping the settlement-pipeline scale moves the \emph{measured} latency ratio across $1$;
Figure~\ref{fig:testbed} overlays the measured $\pwin$ on the closed form. Both exhibit the sharp
threshold at $\rho\approx1$---near-zero below, near-one above---crossing $0.5$ at $\rho\approx1$; the
measured transition is somewhat steeper because the testbed's settlement-latency variance is lower
than the gamma model's. This is corroboration by a separate mechanism, not the model checking itself.
Two concrete observations: the equivocation evidence is a real $\ArtifactEvidenceBytes$-byte object
(two $\ArtifactSigBytes$-byte Ed25519 signatures), and at \emph{realistic} parameters (80\,ms links,
a 2\,s destination block, fast evidence) the measured baseline is $\rho=\ArtifactTbBaseRho$ with
$\pwin=\ArtifactTbBasePwin$ ($\Tsettle\approx\ArtifactTbSettleMs$\,ms vs.\
$\Tacc\approx\ArtifactTbAccMs$\,ms): the naive cross-shard pipeline sits in the \emph{closed-window}
regime, because a multi-second destination block is slower than evidence propagation. The danger is
therefore not committee size but \emph{fast} settlement (a bridge that releases early) against slow or
weak accountability.

\begin{figure}[t]
\centering
\includegraphics[width=\linewidth]{fig_testbed.pdf}
\caption{RQ6. Testbed-\emph{measured} $\pwin$ (points) vs.\ the closed-form model (line) over the
measured latency ratio.}
\label{fig:testbed}
\end{figure}

\paragraph{RQ7: Are deployed systems exposed (documented parameters)?}
Every surveyed deployed system lands in the safe regime---and for a \emph{structural} reason, not an
economic one: it is protected by fast evidence and finality-gating, not by committee size or unbonding
length (Table~\ref{tab:realsystems})~\cite{harmony2019whitepaper,cosmosdocs,ethconsensusspecs,multiversxdocs,polkadotwiki,zilliqadocs}.
\S\ref{sec:measured} then \emph{measures} these load-bearing quantities directly rather than assuming
them; two structural facts explain the placement. First,
equivocation evidence forms in minutes, far inside the unslashable horizon---which \S\ref{sec:measured}
measures as the evidence-expiry window (${\sim}48$\,h on Cosmos), \emph{not} the multi-week unbonding
period it is commonly mistaken for---so $q\approx1$, the bribe is essentially the full stake-at-risk of
$F{+}1$ validators, and the break-even $\Eval^{*}$ sits at $64$--$300$ validator-stakes (G2 hard to
clear). Second, finality-gated settlement (Cosmos IBC, Polkadot relay-chain finality) closes G1 by
construction (Proposition~\ref{prop:immunity}). The only exposed configuration is the fast-exit,
ungated-settlement quadrant---a \emph{hypothetical} row here ($q\approx0.25$, $\Eval^{*}\approx8$
stakes) that \S\ref{sec:measured} then \emph{realizes} on a real value path, where a fast intent
bridge's user leg settles in seconds---the lone quadrant where neither structural protection is present.

\begin{table*}[t]
\centering
\caption{RQ7. The model at documented parameters of named systems (approximate public figures as of
2026). $N$ is the committee size (per-shard committee for sharded rows; validator/finality-committee
size otherwise), distinct from the measurement sample size $n_{\mathrm{obs}}$ of
Table~\ref{tab:measured}. $V^{*}$ is in multiples of one validator's stake-at-risk. \emph{Window (G1)}
is the settlement-gate state (\gated/\openg/\dependsg); \emph{Verdict} is the outcome: \immune\ (G1
closed by construction), \hardgate\ ($q\!\approx\!1$, so G2 is hard to clear), or \exposed. Cosmos~Hub
and Ethereum~PoS are non-sharded accountable-BFT comparators; the fast-exit row is a labelled
\emph{hypothetical}.}
\label{tab:realsystems}
\small
% AUTO-GENERATED by artifact/make_tables.py -- do not edit by hand.
\begin{tabular}{@{}l c r r r r c c@{}}
\toprule
System & Sharded & $N$ & Unbond. & $q$ & $V^{*}$ & Window (G1) & Verdict \\
\midrule
Harmony & yes & 250 & 7\,d & 1.00 & $126$ & \openg & \hardgate \\
MultiversX & yes & 400 & 10\,d & 1.00 & $201$ & \dependsg & \hardgate \\
Zilliqa & yes & 600 & 14\,d & 1.00 & $300$ & \dependsg & \hardgate \\
Cosmos Hub & no & 180 & 21\,d & 1.00 & $90$ & \gated & \immune \\
Ethereum PoS & no & 128 & 1\,d & 1.00 & $64$ & \dependsg & \hardgate \\
Polkadot & yes & 300 & 28\,d & 1.00 & $150$ & \gated & \immune \\
Fast-exit design (hypothetical) & yes & 64 & 0\,d & 0.25 & $8$ & \openg & \exposed \\
\bottomrule
\end{tabular}

\end{table*}

\paragraph{Summary.}
Under the model the attack is profitable only in the joint region
$\{\rho>1\}\cap\{\Eval>\Eval^{*}\}$, the bribe scales only linearly with committee size, and the
defender's decisive lever is the one that closes the window rather than the one that raises the cost.
A discrete-event testbed with real signatures independently reproduces the $\rho\approx1$ threshold,
and a case study of documented parameters places deployed systems in the safe regime---evidence
landing inside the measured evidence-expiry horizon (\S\ref{sec:measured}) keeps $q\approx1$ and (G2)
hard, while finality-gating closes (G1)---so the residual risk is concentrated in fast-exit,
ungated-settlement designs, which \S\ref{sec:measured} exhibits on a real bridge path.
Reconfiguration helps the attacker only in designs that let
old-epoch certified artifacts execute before old-epoch evidence is enforced.
